\chapter{Цифра и аналоговый мир: ЦАП, АЦП, ШИМ}
\label{ch:analog}

В предыдущей главе мы увидели, что ножка микросхемы выдаёт ток, а «0» и «1»~---
диапазоны напряжения. Остался вопрос: а что делать, когда нужен \emph{не} 0 и
не 1? Зуммер должен звучать нужным тоном, светодиод~--- гореть вполсилы, а
датчик температуры выдаёт напряжение, которое надо превратить в число. Эта
глава~--- два моста между цифрой и непрерывным миром.

\begin{buildbox}
Вкладка \emph{Обучение}, книга «из блоков», глава «Подключаем к плате», урок
«Зуммер». Поставьте зуммер на стол Лаборатории и подключите его вход S к
плате. Прочитайте в уроке, чем активный зуммер отличается от пассивного:
какой сигнал нужно подавать на пассивный, чтобы он зазвучал? Запишите
ответ~--- он понадобится в разделе про ШИМ.
\end{buildbox}

Всё, что мы изучали до сих пор, живёт в дискретном мире: сигналы принимают
значения 0 или 1, время идёт тактами. Но окружающий мир непрерывен. Датчик
температуры выдаёт напряжение, микрофон --- переменный ток, а громкоговорителю
нужно именно непрерывное колебание. На границе между цифровой и аналоговой
техникой стоят два типа устройств: \term{АЦП} (аналого-цифровой преобразователь)
превращает напряжение в число, а \term{ЦАП} (цифро-аналоговый преобразователь)
превращает число в напряжение.

\section{Дискретизация и квантование}

Преобразование аналогового сигнала в цифровой состоит из двух шагов. Сначала
напряжение измеряют через равные промежутки времени --- это
\term{дискретизация}, частота измерений называется \term{частотой
дискретизации}. Затем каждое измерение округляют до ближайшего из
$2^N$ уровней --- это \term{квантование}, $N$ --- разрядность преобразователя.

\begin{figure}[H]
\centering
\begin{tikzpicture}[x=1.15cm,y=1.5cm, font=\sffamily\scriptsize]
  \draw[->] (-0.1,0) -- (7.0,0) node[right] {время};
  \draw[->] (0,-1.3) -- (0,1.4) node[above] {$U$};
  \foreach \l in {-1,-0.667,-0.333,0,0.333,0.667,1} \draw[gray!40] (0,\l) -- (6.6,\l);
  \draw[FpgaBlue, thick, domain=0:6.28, samples=120, smooth] plot (\x, {sin(\x r)});
  \foreach \i in {0,...,12} {
    \pgfmathsetmacro{\v}{sin(\i*0.5 r)}
    \pgfmathsetmacro{\q}{round(3*\v)/3}
    \draw[FpgaOrange] (\i*0.5,0) -- (\i*0.5,\q);
    \fill[FpgaOrange] (\i*0.5,\q) circle (1.6pt);
  }
  \node[FpgaBlue, anchor=west] at (4.3,1.15) {непрерывный сигнал};
  \node[FpgaOrange!80!black, anchor=west] at (4.3,0.85) {отсчёты после квантования};
\end{tikzpicture}
\caption{Дискретизация (отсчёты через $0{,}5$ с) и квантование (7 уровней). Оранжевые значения ближе всего к синей кривой, но не совпадают с ней}
\end{figure}

От частоты дискретизации и разрядности зависит качество. Теорема
Котельникова (Найквиста---Шеннона) утверждает: сигнал с наивысшей частотой
$f_{max}$ можно восстановить без потерь, если измерять его с частотой
не менее $2 f_{max}$. Поэтому для звука до 20~кГц диски используют
44,1~кГц. Разрядность определяет шаг квантования: при опорном напряжении
$U_{\text{оп}}$ и разрядности $N$ один \term{младший разряд} (LSB) соответствует
\[
U_{LSB} = \frac{U_{\text{оп}}}{2^N}.
\]
Для 8 бит и 3,3~В получается $3{,}3/256 \approx 12{,}9$~мВ. Каждый
дополнительный бит удваивает число уровней и добавляет около 6~дБ к
отношению сигнал/шум квантования.

\begin{diybox}[Отсчёты и разрядность]
\textbf{Задание.} Увидеть, как таблица отсчётов даёт периодический сигнал, и
оценить влияние разрядности.\\
\textbf{В Logra.} Вкладка \emph{Примеры}: \code{lookup-sine} и
\code{dds-tone}; вкладка \emph{Диаграммы}.\\
\stars{1} \emph{Повтори.} В \code{lookup-sine} сократите шаг и проследите на
диаграмме полный период таблицы отсчётов: сколько отсчётов приходится на
период?\\
\stars{2} \emph{Измени.} Измените один отсчёт таблицы и найдите на диаграмме
такт, на котором форма сигнала изменилась.\\
\stars{3} \emph{Проект.} В \code{dds-tone} задайте слово настройки
$16384$, проследите четыре фазы за период и вычислите частоту сигнала по
формуле $f = f_\text{такт}\cdot F/65536$ (динамик не подключён).\\
\textbf{Ожидаемый результат.} Слово $16384$ даёт четыре отсчёта на период:
частота равна четверти частоты такта.\\
\textbf{Как проверить.} Вкладка \emph{Диаграммы}: период цифрового сигнала
равен четырём тактам.
\end{diybox}

\section{ЦАП: от числа к напряжению}

Самый наглядный ЦАП --- \term{резистивная матрица $R$--$2R$}. Двоичный
код управляет переключателями, каждый разряд вносит в выходное напряжение вклад,
вдвое меньший, чем предыдущий, и на выходе получается число, пропорциональное
коду.

\begin{figure}[H]
\centering
\begin{tikzpicture}[font=\sffamily\scriptsize, res/.style={draw, minimum width=9mm, minimum height=3.6mm, fill=white}]
  \draw (-0.8,0) -- (8.6,0) node[right] {$U_{\text{вых}}$};
  \foreach \i in {0,1,2,3} {
    \pgfmathsetmacro{\x}{\i*2.3+1.15}
    \node[res] at (\i*2.3,0) {$R$};
    \draw (\x,0) -- (\x,-0.6);
    \node[res, rotate=90] at (\x,-1.1) {$2R$};
    \draw (\x,-1.4) -- (\x,-1.9) node[below] {$D_{\i}$};
  }
  \draw (-0.8,0) -- (-0.8,-0.6);
  \node[res, rotate=90] at (-0.8,-1.1) {$2R$};
  \draw (-0.8,-1.4) -- (-0.8,-1.9) node[below] {земля};
  \node[align=left, anchor=west] at (0,1.0) {Разряды $D_0\ldots D_3$ подключаются к $U_{\text{оп}}$ или к земле. Вклад старшего разряда $U_{\text{оп}}/2$, затем $U_{\text{оп}}/4$, \ldots};
\end{tikzpicture}
\caption{Матрица $R$--$2R$ (схема показана упрощённо): каждый разряд управляет переключателем между $U_{\text{оп}}$ и землёй}
\end{figure}

В цифровой микросхеме готового ЦАП обычно нет, но проще всего получить его
из обычных цифровых блоков. Вместо резисторов используют
\term{широтно-импульсную модуляцию} (ШИМ): выход быстро переключается между 0
и 1, причём доля единицы в периоде (скважность) пропорциональна нужному
значению. Простой $RC$-фильтр усредняет эти импульсы и превращает их в
постоянное напряжение. Светодиод, управляемый ШИМ, «дышит», потому что глаз
усредняет импульсы точно так же, как фильтр.

Схема ШИМ состоит из двух частей: восьмиразрядного счётчика, который
циклически считает от 0 до 255, и компаратора, который сравнивает счётчик с
нужным значением \code{level}. Пока счётчик меньше \code{level}, выход равен
единице, иначе нулю. Если, например, \code{level} равно 64, то единица держится
64 такта из 256, то есть среднее напряжение составляет $64/256 = 25\,\%$ от
питания.

\begin{figure}[H]
\centering
\begin{tikzpicture}[x=1.0cm, y=0.9cm, font=\sffamily\small]
  % счётчик: два периода пилы
  \draw[->] (-0.2,0) -- (9.2,0) node[right] {$t$};
  \draw[->] (0,-0.2) -- (0,4.4) node[above] {счётчик};
  \draw[FpgaBlue, very thick] (0,0) -- (4,4) -- (4,0) -- (8,4) -- (8,0);
  \draw[FpgaOrange, thick, dashed] (0,1) -- (9,1) node[right] {\code{level}};
  \node[anchor=east] at (0,4) {255};
  \node[anchor=east] at (0,1) {64};
  % выход
  \begin{scope}[yshift=-3.2cm]
    \draw[->] (-0.2,0) -- (9.2,0) node[right] {$t$};
    \draw[FpgaGreen!60!black, very thick] (0,0) -- (0,1) -- (1,1) -- (1,0) -- (4,0) -- (4,1) -- (5,1) -- (5,0) -- (8,0) -- (8,1) -- (9,1);
    \node[anchor=east] at (0,0.5) {pwm};
    \draw[<->] (0,-0.5) -- node[below] {1 период = 256 тактов} (4,-0.5);
    \draw[<->] (0,1.5) -- node[above] {64 такта} (1,1.5);
  \end{scope}
\end{tikzpicture}
\caption{ШИМ: выход равен единице, пока пила счётчика ниже уровня \code{level}}
\label{fig:pwm_wave}
\end{figure}

Частота ШИМ равна $f_{\text{такт}}/256$: при 27~МГц это около 105~кГц, и
простого $RC$-фильтра с постоянной времени в десятки микросекунд хватает для
сглаживания. Минус метода --- медленная реакция: для нового значения
нужен хотя бы один период.

\subsection{ШИМ на Verilog}

Схема из счётчика и компаратора записывается двумя строками. Первый вариант~---
из файла \code{pwm\_dac.v}: счётчик \code{cnt} бежит по кругу, выход равен
единице, пока \code{cnt < level}.

\vfile{\codepath/00_theory/pwm_dac.v}{code/00\_theory/pwm\_dac.v}

\shot{n-pwm-wave}{ШИМ-ЦАП с кодом 64: выход равен единице первую четверть периода счётчика}{fig:n-pwm}

При \code{level = 64} выход равен единице 64 такта из 256 (рис.~\ref{fig:n-pwm}).
Недостаток простой схемы: если менять \code{level} в середине периода,
импульс получится «рваным». Модуль \code{pwm8} принимает новое значение только
в начале периода: когда счётчик равен 255, значение копируется в регистр
\code{lvl}, и сравнение ведётся с ним.

\vfile{\codepath/b9b_analog/pwm8.v}{code/b9b\_analog/pwm8.v}

\textbf{Как читается.} Два регистра по восемь триггеров: \code{cnt}~--- пила,
\code{lvl}~--- действующий уровень. Выход \code{pwm}~--- компаратор, поэтому
при \code{level = 0} он всегда равен нулю, а при \code{level = 255}~---
255 тактов из 256 (полной единицы такой ШИМ не даёт). Тестбенч
\code{tb\_pwm8.v} для уровней 0, 1, 64, 128, 192, 255 считает единицы за один
полный период и сравнивает с уровнем.

\subsection{Сигма-дельта: единицы вперемешку}

ШИМ собирает все единицы периода в одну пачку. Сигма-дельта модулятор
распределяет их равномерно: накопитель на каждом такте прибавляет уровень, а
\emph{переполнение} накопителя выдаётся как единица на выходе. Остаток от
переполнения остаётся в накопителе и учитывается дальше, поэтому за
$2^W$ тактов единиц будет ровно \code{level}: плотность единиц равна
$\text{level}/2^W$, как у ШИМ, но частота переключений выше, и фильтру
достаточно меньшей постоянной времени.

\vfile{\codepath/b9b_analog/sigma_delta.v}{code/b9b\_analog/sigma\_delta.v}

\textbf{Как читается.} Регистр \code{acc} на один бит шире уровня: старший бит
\code{acc[W]}~--- это перенос прошлого сложения, он и есть выход; младшие
\code{W} бит~--- остаток. При \code{level = 32} (половина шкалы) выход
чередуется 0, 1, 0, 1: единицы стоят через одну. Тестбенч
\code{tb\_sigma\_delta.v} для всех 64 уровней считает единицы в окне из 64
тактов (должно быть ровно \code{level}) и проверяет, что при половинном уровне
единицы не идут подряд.

\begin{diybox}[ШИМ и сигма-дельта]
\textbf{Задание.} Сравнить два способа получить среднее напряжение.\\
\textbf{В Logra.} Вкладка «Код»: \code{pwm8.v} и \code{sigma\_delta.v} (по
одному), «Синтезировать»; «Диаграммы»: число циклов 128, \code{level = 16} у
сигма-дельты (\code{W = 6}) и \code{level = 64} у ШИМ; вкладка «Примеры»:
\code{pwm}, \code{sigma-delta}.\\
\textbf{Ожидаемый результат.} В ШИМ единицы идут пачкой, у сигма-дельты они
разбросаны; в окне из 64 тактов у сигма-дельты ровно 16 единиц.\\
\textbf{Как проверить.} Тестбенчи печатают \code{pwm8: OK} и
\code{sigma\_delta: OK}. \stars{2} Допишите ШИМ на 4 бита (16 уровней).
\end{diybox}

\textbf{Образец ответа.}
\begin{vcode}
module pwm4 (
    input        clk,
    input  [3:0] level,
    output       pwm
);
    reg [3:0] cnt = 4'd0;
    always @(posedge clk) cnt <= cnt + 4'd1;
    assign pwm = (cnt < level);
endmodule
\end{vcode}


\subsection{Звук: зуммер на ШИМ}

Пассивный зуммер воспроизводит ту частоту, которую ему подают. Поэтому для
звука нужен не уровень, а \emph{периодический сигнал}: прямоугольные
импульсы с частотой тона. Чтобы получить, например, тон 2~кГц от тактового
сигнала 27~МГц, надо переключать выход каждые $27\,000\,000/(2\cdot 2000)
= 6750$ тактов. Скважность задаёт громкость: при 50\,\% звук самый громкий, при
малой доле единицы~--- тише.

Схема зуммера на Verilog~--- счётчик, который считает до \code{PERIOD-1} и
сбрасывается, и компаратор: выход равен единице, пока счётчик меньше
\code{HIGH}. Разрешение \code{en} (в проекте это состояние ТРЕВОГА) обнуляет
счётчик и выход.

\vfile{\codepath/b9b_analog/buzzer.v}{code/b9b\_analog/buzzer.v}

\textbf{Как читается.} По умолчанию \code{PERIOD = 13500} и \code{HIGH = 3375}:
при 27~МГц это тон 2~кГц со скважностью 25\,\%. Выход \code{out} не
запоминается в триггере, а вычисляется из \code{en} и счётчика, поэтому после
снятия тревоги звук прекращается сразу. Тестбенч \code{tb\_buzzer.v} за два
периода считает 6750 единиц, измеряет расстояние между фронтами (13\,500
тактов) и убеждается, что при \code{en = 0} выход молчит.

\begin{diybox}[Этап 10 проекта: зуммер на ШИМ]
\textbf{Задание.} Добавить к панели звуковой сигнал тревоги и подключить его к
плате.\\
\textbf{В Logra.} Вкладка \emph{Обучение}: урок «Зуммер» (глава «Подключаем к
плате») и урок «Кнопка на настоящей плате»; вкладка \emph{Примеры}: \code{pwm},
\code{breathing-pwm}; вкладки \emph{Диаграммы}, \emph{Выводы}, \emph{Подключение}.\\
\stars{1} \emph{Повтори.} Урок «Зуммер»: поставьте зуммер на стол и подключите
его вход S к плате. В примере \code{pwm} проследите, как уровень задаёт долю
единицы.\\
\stars{2} \emph{Измени.} Задача~4: подайте код 192 и измерьте скважность на
\emph{Диаграммах} (ожидаемо 75\,\%). Затем в \code{breathing-pwm} сократите
\code{LEVEL\_TICKS} и найдите развороты уровня на границах 0 и 255.\\
\stars{3} \emph{Проект.} \textbf{Этап проекта «Охранная панель»:} зуммер
на ШИМ. Состояние ТРЕВОГА (этап~9) включает тон 2~кГц со скважностью 25\,\%;
вне тревоги выход зуммера равен нулю. Критерий приёмки: на диаграмме
период тона равен 13\,500 тактов при тактовой частоте 27~МГц, доля единицы
в периоде равна 25\,\%, а при выходе из тревоги звук прекращается на
следующем такте. Затем подключите зуммер на вкладке \emph{Подключение}
(назначьте порт во вкладке \emph{Выводы}), прошейте и убедитесь, что тон
слышен.\\
\textbf{Ожидаемый результат.} Зуммер звучит ровно в состоянии ТРЕВОГА и
молчит в остальных.\\
\textbf{Как проверить.} Вкладка \emph{Диаграммы}: выход зуммера и состояние
панели на одной шкале времени; вкладка \emph{Выводы}: порт зуммера назначен.\\
\textbf{То же на Verilog.} Вкладка \emph{Код}: вставьте \code{buzzer.v},
«Синтезировать»; на «Диаграммах» подайте \code{en = 1} и посчитайте период.
Чтобы увидеть период на экране, при создании экземпляра задайте
\code{buzzer \#(.PERIOD(20), .HIGH(5))}. Образец ответа к \stars{3}: сигнал
ТРЕВОГА из автомата панели подключают к \code{en}, выход~--- к порту зуммера.
\end{diybox}

\begin{vcode}
module alarm_sound (
    input  clk,
    input  alarm,                // 1 в состоянии ТРЕВОГА
    output buzzer_pin
);
    reg [13:0] cnt = 14'd0;
    always @(posedge clk) begin
        if (!alarm || cnt == 14'd13499) cnt <= 14'd0;
        else                            cnt <= cnt + 14'd1;
    end
    assign buzzer_pin = alarm && (cnt < 14'd3375);   // 25 % от 13500
endmodule
\end{vcode}


\section{АЦП: от напряжения к числу}

Существует несколько способов построить АЦП, и у каждого есть своя цена.

\begin{table}[H]
\centering
\caption{Основные типы АЦП}
\small
\begin{tabularx}{\textwidth}{l X X}
\toprule
\textbf{Тип} & \textbf{Как работает} & \textbf{Особенность} \\
\midrule
Параллельный (flash) & $2^N-1$ компараторов сравнивают вход с лесенкой напряжений одновременно & Мгновенный, но огромный: для 8 бит нужно 255 компараторов \\
Последовательного приближения (SAR) & $N$ тактов поиска по методу бисекции с одним компаратором и ЦАП & Золотая середина: быстрый и дешёвый, самый распространённый \\
Сигма-дельта & Однобитный АЦП с очень высокой частотой и цифровым фильтром & Очень точный, но медленный; применяется в звуке и измерителях \\
Интегрирующий & Измеряет время разряда конденсатора & Очень точный, но медленный; мультиметры \\
\bottomrule
\end{tabularx}
\end{table}

\subsection{АЦП последовательного приближения}

Идея похожа на игру «угадай число от 0 до 255»: каждый раз называем середину
оставшегося интервала и по ответу «больше» или «меньше» отбрасываем половину.
За $N$ вопросов находим число с точностью до одного разряда. Роль
«вопроса» играет компаратор, который сравнивает входное напряжение с
напряжением ЦАП, а роль ведущего --- цифровой автомат, который по очереди
проверяет биты от старшего к младшему.

\begin{figure}[H]
\centering
\begin{tikzpicture}[font=\sffamily\small, bx/.style={draw, minimum width=26mm, minimum height=11mm, fill=white, rounded corners=2pt, align=center}]
  \node[bx] (cmp) at (0,0) {Компаратор};
  \node[bx, minimum width=36mm] (sar) at (5.2,0) {Автомат SAR\\{\scriptsize регистр + управление}};
  \node[bx] (dac) at (5.2,-2.2) {ЦАП};
  \draw[-{Stealth}] (-2.6,0.25) node[left] {$U_{\text{вх}}$} -- ([yshift=2.5mm]cmp.west);
  \draw[-{Stealth}] (dac.west) -| node[pos=0.25, above] {$U_{\text{ЦАП}}$} ([xshift=-4mm]cmp.south);
  \draw[-{Stealth}] (cmp) -- node[above] {cmp} (sar);
  \draw[-{Stealth}] (sar) -- node[right] {код} (dac);
  \draw[-{Stealth}] (sar.east) -- ++(1.4,0) node[right] {результат};
\end{tikzpicture}
\caption{Структура АЦП последовательного приближения. Цифровая часть --- конечный автомат, а значит, она целиком реализуется цифровой схемой}
\end{figure}

Цифровая часть --- конечный автомат с регистром результата: он по сигналу
\code{cmp} решает, оставить ли проверяемый бит единицей. Проследим его работу для
входного напряжения, которое соответствует числу 179 ($\text{B3}_{16}=10110011_2$).

\begin{table}[H]
\centering
\caption{Поиск числа 179 методом последовательного приближения}
\small
\begin{tabular}{clccl}
\toprule
\textbf{Шаг} & \textbf{Пробуем (код ЦАП)} & \textbf{В десятичной} & \textbf{$U_{\text{вх}}\ge U_{\text{ЦАП}}$?} & \textbf{Результат после шага} \\
\midrule
1 & \code{0x80} = \code{10000000} & 128 & да & \code{0x80} \\
2 & \code{0xC0} = \code{11000000} & 192 & нет & \code{0x80} \\
3 & \code{0xA0} = \code{10100000} & 160 & да & \code{0xA0} \\
4 & \code{0xB0} = \code{10110000} & 176 & да & \code{0xB0} \\
5 & \code{0xB8} = \code{10111000} & 184 & нет & \code{0xB0} \\
6 & \code{0xB4} = \code{10110100} & 180 & нет & \code{0xB0} \\
7 & \code{0xB2} = \code{10110010} & 178 & да & \code{0xB2} \\
8 & \code{0xB3} = \code{10110011} & 179 & да & \code{0xB3} \\
\bottomrule
\end{tabular}
\end{table}

Восемь шагов дают восемь бит результата: на каждом шаге автомат проверяет
ровно один бит, начиная со старшего.

\subsection{АЦП последовательного приближения на Verilog}

Автомат записывается в одном модуле \code{sar\_adc}: регистр \code{acc} хранит
уже найденные старшие биты, «бегущая единица» \code{mask} указывает проверяемый
бит, а сигнал \code{cmp} от внешнего компаратора решает, остаётся ли этот бит
единицей. В файле есть и учебная обёртка \code{sar\_demo}: «входное напряжение»
в ней задано константой 179, а компаратор собран из оператора сравнения, так что
всю систему можно проверить без внешних деталей.

\vfile{\codepath/00_theory/sar_adc.v}{code/00\_theory/sar\_adc.v}

\textbf{Как читается.} Выражение \code{acc\_n = cmp ? (acc | mask) : acc}~---
это решение на текущем шаге: если вход не меньше пробного значения, бит
остаётся единицей. Когда \code{mask[0]} равно единице, проверен самый младший
бит: автомат записывает \code{result}, поднимает \code{done} на один такт и
возвращается в ожидание. Иначе \code{mask} сдвигается вправо, а \code{dac}
получает следующую пробу.

\shot{n-sar-wave}{Поиск числа 179: код ЦАП проходит значения \code{80}, \code{C0}, \code{A0}, \code{B0}, \code{B8}, \code{B4}, \code{B2}, \code{B3}}{fig:n-sar}

На диаграмме (рис.~\ref{fig:n-sar}) виден тот же бинарный поиск, что и в
таблице: через восемь тактов \code{result} становится равным \code{0xB3}, то
есть 179. Тестбенч \code{00\_theory/tb\_sar\_adc.v} перебирает входные
значения с шагом 7 от 0 до 255 и сравнивает \code{result} с входом для каждого;
он печатает \code{sar\_adc: OK}.

\begin{infobox}[АЦП и цифровая микросхема]
Во многих цифровых микросхемах «медленного» АЦП, как у микроконтроллера, нет.
Если нужно измерять напряжение, подключают внешнюю микросхему АЦП по
последовательному интерфейсу. \lookahead{на плате это именно так;
интерфейсы разобраны в главе~\ref{ch:uart}.}
\end{infobox}

\begin{diybox}[АЦП: поиск числа по битам]
\textbf{Задание.} Проследить работу АЦП последовательного приближения и
увидеть, что результат даёт цифровой автомат, а не аналоговая схема.\\
\textbf{В Logra.} Вкладка \emph{Примеры}: \code{sigma-delta};
вкладки \emph{Диаграммы}, \emph{RTL}.\\
\stars{1} \emph{Повтори.} Проведите вручную поиск числа 179 по таблице главы
и выпишите пробные коды ЦАП.\\
\stars{2} \emph{Измени.} Задача~6: выпишите пробные коды для числа 100 и
сравните с вариантом ответа.\\
\stars{3} \emph{Проект.} В примере \code{sigma-delta} после синхронизации
уровня посчитайте 64 такта: уровень 16 даёт 16 единиц, уровень 48~--- 48.
Объясните, чем плотность единиц похожа на ШИМ и чем отличается.\\
\textbf{Ожидаемый результат.} Количество единиц за 64 такта равно уровню.\\
\textbf{Как проверить.} Вкладка \emph{Диаграммы}: подсчёт единиц за 64 такта;
вкладка \emph{RTL}: найдите аккумулятор (регистр остатка) и его перенос.\\
\textbf{То же на Verilog.} Вкладка \emph{Код}: вставьте \code{sar\_adc.v}, верхним
выберите \code{sar\_demo}, «Синтезировать»; «Диаграммы»: число циклов 20,
импульс \code{start}; сверьте значения \code{dac} с таблицей главы.
\stars{2} Измерьте число 100. \textbf{Образец ответа:} в строке
\code{localparam [7:0] VIN = 8'd179;} замените 179 на 100; пробные коды
\code{0x80, 0x40, 0x60, 0x70, 0x68, 0x64, 0x66, 0x65}, результат \code{0x64}.
Тестбенч \code{00\_theory/tb\_sar\_adc.v} проверяет сразу 37 входных значений.
\end{diybox}

\begin{warnbox}[Типичные ошибки]
\begin{itemize}
  \item \textbf{Путать разрядность и точность.} Много бит даёт мелкий шаг
        $U_\text{LSB}$, но не уберёт шум и погрешность опорного напряжения.
  \item \textbf{Брать частоту дискретизации ровно $2f_\text{max}$ без запаса.}
        Реальные сигналы шире идеальных; частоту берут с запасом.
  \item \textbf{Подавать ШИМ без фильтра туда, где нужно постоянное
        напряжение.} Без фильтрации получатся не «среднее напряжение», а
        быстрые импульсы.
  \item \textbf{Ждать мгновенной реакции ШИМ-ЦАП.} Новое значение полностью
        появляется только через период, поэтому медленный ШИМ нельзя
        использовать для быстрого сигнала.
\end{itemize}
\end{warnbox}

\section{Вопросы и задачи}
\label{sec:tasks-analog}

Ответы собраны в приложении~\ref{app:answers}.

\begin{taskbox}[Задачи к главе]
\begin{enumerate}
  \item \stars{1}~10-разрядный АЦП с опорным напряжением 3,3~В. Какому изменению
        напряжения соответствует один младший разряд?
  \item \stars{1}~Сколько бит нужно, чтобы различать напряжения с шагом не больше
        10~мВ в диапазоне 0--5~В?
  \item \stars{1}~С какой минимальной частотой надо измерять сигнал с частотой
        до 3~кГц, чтобы по отсчётам его можно было восстановить?
  \item \stars{1}~На вход ШИМ-ЦАП подали код 192. Какова доля времени, в которую
        выход равен единице? Какое среднее напряжение получится при питании
        3,3~В?
  \item \stars{2}~Сколько тактов преобразования нужно 12-разрядному АЦП
        последовательного приближения? Сколько компараторов нужно параллельному
        АЦП на 12 бит?
  \item \stars{2}~Опишите, какие значения будет перебирать АЦП последовательного
        приближения при измерении кода 100 (восемь бит).
  \item \stars{2}~\emph{Найдите ошибку.} В поиске числа 100 студент записал:
        «шаг 1: пробуем \code{0x80}, ответ нет, результат \code{0x80}; шаг 2:
        пробуем \code{0xC0}\ldots». Что неверно и какая пробная
        комбинация должна быть на втором шаге?
  \item \stars{3}~\emph{Мини-проект.} Постройте восьмиразрядный ШИМ из счётчика и
        компаратора. Критерий приёмки: на диаграмме при \code{level}=64 выход
        равен единице 64 такта из 256, при \code{level}=0 выход всё время
        равен нулю, при \code{level}=255 --- 255 тактов из 256.
  \item \stars{3}~\emph{Мини-проект (этап~10).} Выход зуммера панели: в состоянии
        ТРЕВОГА он выдаёт тон 2~кГц со скважностью 25\,\% (тактовая частота
        27~МГц), иначе равен нулю. Критерий приёмки: на диаграмме период равен
        13\,500 тактов, единица держится 3375 тактов, при выходе из тревоги
        звук прекращается.
\end{enumerate}
\end{taskbox}

\section{Что дальше}

Мы прошли путь от сигнала до памяти, автоматов, ножки микросхемы и связи с
аналоговым миром. Этого достаточно, чтобы понять, как устроено практически
любое цифровое устройство: процессор, контроллер дрона, видеокарта. К этому
моменту у вас собрана «Охранная панель»: датчик, логика тревоги, запоминание,
счётчик времени, автомат режимов с вводом кода, индикатор и зуммер.

Остаётся связать теорию с устройством ПЛИС и маршрутом до платы, которые мы уже
разобрали: как каждый узел, собранный из блоков, записывается кодом на Verilog и
какому ресурсу микросхемы он соответствует? Это тема главы~\ref{ch:bridge}.
Ниже~--- «карта перехода»: какому ресурсу микросхемы соответствует каждый узел и
в каком задании Logra вы его уже собирали.

{\small
\begin{longtable}{@{}>{\raggedright\arraybackslash}p{0.30\linewidth}>{\raggedright\arraybackslash}p{0.30\linewidth}>{\raggedright\arraybackslash}p{0.32\linewidth}@{}}
\caption*{\textbf{\sffamily Карта перехода от схем к ПЛИС}}\\
\toprule
\textbf{Узел} & \textbf{Ресурс ПЛИС} & \textbf{Задание в Logra} \\
\midrule
\endfirsthead
\toprule
\textbf{Узел} & \textbf{Ресурс ПЛИС} & \textbf{Задание в Logra} \\
\midrule
\endhead
Логические элементы, дешифраторы, шифраторы, мультиплексоры & таблицы LUT & урок «Мультиплексор 4→1»; пример \code{mux-four} \\
Сумматоры, счётчики & LUT + быстрые цепочки переноса (ALU) & урок «Полный сумматор»; урок «Счётчик одним блоком» \\
Умножители & аппаратные блоки DSP & пример \code{shift-add-multiply}; пример \code{alu} \\
Триггеры, регистры, сдвиговые регистры & триггеры в каждой логической ячейке & урок «D-триггер: запомнить бит»; пример \code{shift} \\
Конечные автоматы & LUT + триггеры & урок «Кодовый замок»; пример \code{door-controller} \\
ПЗУ, ОЗУ и FIFO & блочная память BSRAM, распределённая память SSRAM & примеры \code{small-rom}, \code{ram-write-read}, \code{fifo-four} \\
Большая память & встроенная SDRAM 64 Мбит & пример \code{ram-write-read} в малом масштабе \\
Тактовый сигнал & генератор 27 МГц, PLL, глобальная тактовая сеть & урок «Делитель частоты на 2»; пример \code{blink} \\
Выходы с тремя состояниями и подтяжки & блоки ввода-вывода IOBUF & вкладка \emph{Выводы}; урок «Кнопка управляет светодиодом» \\
ШИМ, ЦАП и АЦП & счётчик и компаратор + внешний фильтр; внешняя микросхема АЦП & примеры \code{pwm}, \code{breathing-pwm}, \code{sigma-delta}; урок «Зуммер» \\
\bottomrule
\end{longtable}}

\begin{ideabox}
Комбинационная логика вычисляет, последовательная запоминает. Синхронная
схема --- это регистры и комбинационная логика между ними с одним общим тактом.
Её быстродействие определяется критическим путём. Конечный автомат ---
регистр состояния плюс логика переходов. Память --- массив ячеек с дешифратором
адреса. На границе кристалла работают токи, пороги и подтяжки, а аналоговый
мир подключается через ЦАП и АЦП. В следующих главах все эти узлы мы
запишем на Verilog и соберём на настоящей плате.
\end{ideabox}

\begin{buildbox}[Первый шаг к плате]
Откройте вкладку \emph{Обучение}, урок «Кнопка на настоящей плате». Это самый
короткий путь от схемы к устройству: подключите плату, назначьте порты на
вкладке \emph{Выводы} и убедитесь, что кнопка зажигает светодиод. Затем
вернитесь к проекту «Охранная панель» и соберите его целиком на плате.
Подключение платы разобрано в главе~\ref{ch:board}, первый проект~---
в главе~\ref{ch:examples}.
\end{buildbox}
